\documentclass[10pt,a4paper]{amsart}
\usepackage[margin=19mm]{geometry}
\usepackage{amsmath,amssymb,mathtools}
\usepackage[scr=rsfs,cal=euler]{mathalfa}
\usepackage{xfrac}
\usepackage[unicode,hidelinks,bookmarks=false]{hyperref}
\newcommand{\ie}{i.e.\ }
\newcommand{\cf}{cf.\ }
\newcommand{\resp}{respectively\ }
\newcommand{\Subsection}{\subsection}
\providecommand{\nobreakdash}{\nobreak\hbox{-}}
\setlength{\emergencystretch}{2em}
\multlinegap0pt
\usepackage{amssymb}
\usepackage{mathtools}
\usepackage{xcolor}
\usepackage[shortlabels]{enumitem}
\newtheorem{proposition}{Proposition}
\title{From Linear to Nonlinear: A Resolvente criterion for Polynomial Stability of Semigroups Generated by Monotone Operators}
\author{Marcelo M. Cavalcanti, Val\'eria N. Domingos Cavalcanti, Jaime E. Mun\~oz Rivera}
\date{Source: arXiv:2604.03811v1 (CC BY 4.0)}
\begin{document}
\maketitle

\noindent\emph{Source and licence.} From Linear to Nonlinear: A Resolvente criterion for Polynomial Stability of Semigroups Generated by Monotone Operators. Version arXiv:2604.03811v1 (CC BY 4.0), distributed by arXiv under the Creative Commons Attribution 4.0 International licence, http://creativecommons.org/licenses/by/4.0/, as recorded in the arXiv licence metadata for this exact version and shown on the abstract page https://arxiv.org/abs/2604.03811v1. Full licence notice: \texttt{licenses/arxiv-2604.03811-CC-BY-4.0.txt}.\par\smallskip

\noindent\textit{Editorial scope.} Counted unit: the article's proposition prop:resolvent\_app3 (source0.tex lines 1048--1054). The article's complete original proof of it is delivered with it (source0.tex lines 1056--1100, one \texttt{proof} environment of 45 lines). No mathematics is written, repaired or re-derived: the statement and the proof are the article's own lines, in the article's own notation, and no retained line is substituted or re-flowed.  No further source-local context is needed: every cross-reference of every retained line resolves inside the two delivered spans.  Nothing was omitted from any retained span, so the package carries no omission record.  No retained line contains a citation, so the wrapper carries no bibliography and no bibliography entry.  No retained line contains a figure environment, an image-inclusion command or drawing code, so no image is delivered and the package declares no asset dependency.  Editorial formatting only, in addition: an AMS \texttt{amsart} preamble that restates the article's own declarations and macros used by the retained lines, namely the environments \texttt{proposition} on separate counters, together with the standard packages those lines need.  The retained environments print in this document as Proposition 1 in order of appearance, whereas the article prints the same items with the numbering of its own full document.  The complete native source of the article as distributed in the arXiv e-print for this exact version is bundled unchanged as \texttt{sources/197835/source0.tex}.
\begin{proposition}\label{prop:resolvent_app3}
Let $U_\lambda=(u_\lambda, v_\lambda)$ be the solution to $(\lambda I + \mathcal{A})U_\lambda = F$ for $F=(f,g) \in \mathcal{H}$. Then, as $\lambda \to 0^+$, the solution satisfies the growth estimate:
\[
\|U_\lambda\|_{\mathcal H} \le \frac{C}{\lambda} \|F\|_{\mathcal H},
\]
where $C > 0$ depends on the operator $A$ but is independent of $\lambda$.
\end{proposition}

\begin{proof}
The resolvent system reads:
\begin{align}
\lambda u_\lambda - v_\lambda &= f, \label{eq:res_nl_1} \\
\lambda v_\lambda + A u_\lambda + \rho(\|\mathcal C v_\lambda\|_H^2) v_\lambda &= g. \label{eq:res_nl_2}
\end{align}
Taking the inner product of the second equation \eqref{eq:res_nl_2} with $v_\lambda$ in $H$, we obtain:
\begin{equation}\label{eq:resolvent_test_app3}
\lambda \|v_\lambda\|_H^2 + \langle A u_\lambda, v_\lambda \rangle_H + \rho(\|\mathcal C v_\lambda\|_H^2) \|v_\lambda\|_H^2 = \langle g, v_\lambda \rangle_H.
\end{equation}
From the first equation \eqref{eq:res_nl_1}, we have $u_\lambda = \frac{v_\lambda + f}{\lambda}$. Since $A$ is a positive self-adjoint operator, we can rewrite the cross-term using fractional powers:
\begin{align*}
\langle A u_\lambda, v_\lambda \rangle_H &= \langle A^{1/2} u_\lambda, A^{1/2} v_\lambda \rangle_H \\
&= \left\langle A^{1/2} \left(\frac{v_\lambda + f}{\lambda}\right), A^{1/2} v_\lambda \right\rangle_H \\
&= \frac{1}{\lambda} \|A^{1/2} v_\lambda\|_H^2 + \frac{1}{\lambda} \langle A^{1/2} f, A^{1/2} v_\lambda \rangle_H.
\end{align*}
Substituting this back into \eqref{eq:resolvent_test_app3}, we get the fundamental identity:
\begin{equation}
\lambda \|v_\lambda\|_H^2 + \frac{1}{\lambda} \|A^{1/2} v_\lambda\|_H^2 + \rho(\|\mathcal C v_\lambda\|_H^2) \|v_\lambda\|_H^2 = \langle g, v_\lambda \rangle_H - \frac{1}{\lambda} \langle A^{1/2} f, A^{1/2} v_\lambda \rangle_H.
\end{equation}

As $\lambda \to 0^+$, the singular term $\frac{1}{\lambda} \|A^{1/2} v_\lambda\|_H^2$ provides the dominant coercivity. Discarding the non-negative terms $\lambda \|v_\lambda\|_H^2$ and $\rho(\|\mathcal C v_\lambda\|_H^2) \|v_\lambda\|_H^2$ on the left-hand side, and applying the Cauchy-Schwarz inequality, we have:
\begin{align*}
\frac{1}{\lambda} \|A^{1/2} v_\lambda\|_H^2 &\le \|g\|_H \|v_\lambda\|_H + \frac{1}{\lambda} \|A^{1/2} f\|_H \|A^{1/2} v_\lambda\|_H.
\end{align*}
Since $A$ has a compact inverse (Assumption H1), there exists a Poincaré-type constant $C_P > 0$ such that $\|v_\lambda\|_H \le C_P \|A^{1/2} v_\lambda\|_H$. Thus:
\begin{equation}
\frac{1}{\lambda} \|A^{1/2} v_\lambda\|_H^2 \le C_P \|g\|_H \|A^{1/2} v_\lambda\|_H + \frac{1}{\lambda} \|A^{1/2} f\|_H \|A^{1/2} v_\lambda\|_H.
\end{equation}
Dividing by $\frac{1}{\lambda} \|A^{1/2} v_\lambda\|_H$ (assuming $v_\lambda \neq 0$), we conclude:
\begin{equation}
\|A^{1/2} v_\lambda\|_H \le C_P \lambda \|g\|_H + \|A^{1/2} f\|_H.
\end{equation}
This crucial estimate shows that the "elastic" energy of the velocity, $\|A^{1/2} v_\lambda\|_H$, remains uniformly bounded as $\lambda \to 0^+$.

Finally, to estimate the displacement $u_\lambda$, we use $u_\lambda = \frac{v_\lambda + f}{\lambda}$:
\begin{align*}
\|A^{1/2} u_\lambda\|_H &\le \frac{1}{\lambda} \left( \|A^{1/2} v_\lambda\|_H + \|A^{1/2} f\|_H \right) \\
&\le \frac{1}{\lambda} \left( C_P \lambda \|g\|_H + 2\|A^{1/2} f\|_H \right).
\end{align*}
Since $\|v_\lambda\|_H \le C_P \|A^{1/2} v_\lambda\|_H$, it follows that the full norm $\|U_\lambda\|_{\mathcal{H}}^2 = \|A^{1/2} u_\lambda\|_H^2 + \|v_\lambda\|_H^2$ behaves as $\mathcal{O}(\lambda^{-1})$:
\begin{equation}
\|U_\lambda\|_{\mathcal H} \le \frac{C}{\lambda} \|F\|_{\mathcal H}, \qquad \text{as } \lambda \to 0^+.
\end{equation}
\end{proof}

\end{document}
